\documentclass[11pt]{article}

\usepackage[margin=1in]{geometry}
\usepackage{amsmath,amssymb}
\usepackage{hyperref}
\usepackage{enumitem}

\hypersetup{
  colorlinks=true,
  linkcolor=blue,
  urlcolor=blue,
  citecolor=blue
}
\emergencystretch=2em

\newcommand{\Ball}{B}
\newcommand{\ellone}{\ell_1}
\newcommand{\ellinf}{\ell_\infty}

\title{A Literature-Implied Negative Answer to the Cotype-Only Helly Question}
\author{Math discovery packet}
\date{June 16, 2026}

\begin{document}
\maketitle

\section{Status}

This is a lightweight \emph{literature-implied answer}, not an original
solution.  The source paper is
\[
\text{G. Ivanov, \emph{No-dimensional Helly's theorem in uniformly convex
Banach spaces}, arXiv:2409.05744.}
\]
In the Introduction, immediately before Section 2, the paper asks:
\begin{quote}
Is there a version of the no-dimensional Helly's theorem and a corresponding
dual version of Maurey's lemma in a Banach space of cotype \(q\ge 2\)?
\end{quote}

The supporting paper is
\[
\text{G. Ivanov, \emph{On Banach Spaces with the Helly Approximation Property},
arXiv:2603.22743.}
\]

\section{Identification}

ArXiv:2603.22743 defines a Helly approximation sequence exactly in the
``no-dimensional Helly'' sense: every finite family of convex sets whose
\(k\)-subfamilies meet the unit ball has a common point in the
\(r_k\)-neighborhoods of all members.  The space has the Helly approximation
property when one can choose \(r_k\to0\).

The supporting paper proves the following two facts.
\begin{enumerate}[leftmargin=*]
\item Theorem 1.1 gives a no-dimensional Helly theorem when \(X^*\) has
Rademacher type \(p>1\), with
\[
  r_k(X)=6T_p(X^*)k^{-1+1/p}.
\]
\item Theorem 1.3 characterizes the property: \(X\) has the Helly
approximation property if and only if \(X^*\), equivalently \(X\), has
non-trivial Rademacher type.
\end{enumerate}

Consequently, the correct qualitative hypothesis is non-trivial type, not
cotype alone.  Indeed, \(L_1[0,1]\) has cotype \(2\), while it has only
trivial type because it contains \(\ell_1^n\)'s uniformly, for example on
disjoint measurable sets.  By the characterization in arXiv:2603.22743,
\(L_1[0,1]\) does not have the Helly approximation property.  Therefore no
cotype-only theorem can guarantee a no-dimensional Helly radius tending to
zero.

\section{Interpretation}

This also explains the failure of the requested ``dual Maurey's lemma'' under
cotype alone.  The dual reduction in arXiv:2603.22743 turns Helly approximation
into an averaging statement in \(X^*\).  For \(X=L_1[0,1]\), the dual is
\(L_\infty[0,1]\), which has no non-trivial Rademacher type, so Maurey's
sampling mechanism is unavailable.

\section{Scope}

The supporting theorem does not rule out cotype appearing in sharper
quantitative statements once non-trivial type is also present.  It rules out
the source question as a cotype-only principle: finite cotype, even cotype
\(2\), is insufficient.

\section{References}

\begin{itemize}[leftmargin=*]
\item G. Ivanov, \emph{No-dimensional Helly's theorem in uniformly convex
Banach spaces}, arXiv:2409.05744.
\item G. Ivanov, \emph{On Banach Spaces with the Helly Approximation Property},
arXiv:2603.22743.
\item J. Lindenstrauss and L. Tzafriri, \emph{Classical Banach Spaces II:
Function Spaces}, for the standard type and cotype facts used for \(L_1\).
\end{itemize}

\end{document}
